% from ega2-1-fr.tex
\oldpage[II]{12}
\begin{proposition}[1.4.11]
\label{II.1.4.11-fr}
Soient $S$ un préschéma, $\mathcal{B}$ une $\mathcal{O}_S$-Algèbre
quasi-cohérente, $X=\operatorname{Spec}(\mathcal{B})$. Pour tout Idéal
quasi-cohérent $\mathcal{K}$ de $\mathcal{O}_S$, on a (en désignant par $f$
le morphisme structural $X\to S$),
$f^*(\mathcal{K})\mathcal{O}_X=\widetilde{\mathcal{K}\mathcal{B}}$ à un
isomorphisme canonique près.
\end{proposition}

La question étant locale sur $S$, on est ramené au cas où
$S=\operatorname{Spec}(A)$ est affine, et dans ce cas la proposition n'est
autre que (\hyperref[I.1.6.9-fr]{I, 1.6.9}).

\subsection{Changement du préschéma de base}
\label{II.1.5-fr}

\begin{proposition}[1.5.1]
\label{II.1.5.1-fr}
Soit $X$ un préschéma affine au-dessus de $S$. Pour toute extension
$g:S'\to S$ du préschéma de base, $X'=X_{(S')}=X\times_S S'$ est affine
au-dessus de $S'$.
\end{proposition}

Si $f'$ est la projection $X'\to S'$, il suffit de prouver que
$f^{\prime-1}(U')$ est un ouvert affine pour tout ouvert affine $U'$ de $S'$
tel que $g(U')$ soit contenu dans un ouvert affine $U$ de $S$
(\hyperref[II.1.2.1-fr]{1.2.1})~; on peut donc se limiter au cas où $S$ et
$S'$ sont affines, et il suffit de prouver que $X'$ est alors un schéma
affine (\hyperref[II.1.3.4-fr]{1.3.4}). Mais alors
(\hyperref[II.1.3.4-fr]{1.3.4}) $X$ est un schéma affine, et si $A$, $A'$ et
$B$ sont les anneaux de $S$, $S'$ et $X$ respectivement, on sait que $X'$ est
un schéma affine d'anneau $A'\otimes_A B$
(\hyperref[I.3.2.2-fr]{I, 3.2.2}).

\begin{corollary}[1.5.2]
\label{II.1.5.2-fr}
Sous les hypothèses de (\hyperref[II.1.5.1-fr]{1.5.1}), soient $f:X\to S$
le morphisme structural, $f':X'\to S'$, $g':X'\to X$ les projections, de
sorte que le diagramme
\[
  \xymatrix{
    X\ar[d] & X'\ar[l]_{g'}\ar[d]^{f'}\\
    S & S'\ar[l]_g
  }
\]
est commutatif. Pour tout $\mathcal{O}_X$-Module quasi-cohérent $\mathcal{F}$,
il existe un isomorphisme canonique de $\mathcal{O}_{S'}$-Modules
\[
  u:g^*(f_*(\mathcal{F}))\xrightarrow{\sim}
  f'_*(g^{\prime*}(\mathcal{F})).
  \tag{1.5.2.1}\label{II.1.5.2.1-fr}
\]
En particulier, il existe un isomorphisme canonique de $\mathcal{A}(X')$ sur
$g^*(\mathcal{A}(X))$.
\end{corollary}

Pour définir $u$, il suffit de définir un homomorphisme
\[
  v:f_*(\mathcal{F})\longrightarrow
  g_*(f'_*(g^{\prime*}(\mathcal{F})))
  =f_*(g'_*(g^{\prime*}(\mathcal{F})))
\]
et de prendre $u=v^\sharp$ (\hyperref[0.4.4.3-fr]{0, 4.4.3}). Nous prendrons
$v=f_*(\rho)$, où $\rho$ est l'homomorphisme canonique
$\mathcal{F}\to g'_*(g^{\prime*}(\mathcal{F}))$
(\hyperref[0.4.4.3-fr]{0, 4.4.3}). Pour prouver que $u$ est un isomorphisme,
on peut se borner au cas où $S$ et $S'$, donc $X$ et $X'$, sont affines~;
avec les notations de (\hyperref[II.1.5.1-fr]{1.5.1}), on a alors
$\mathcal{F}=\widetilde{M}$, où $M$ est un $B$-module. On constate alors
aussitôt que $g^*(f_*(\mathcal{F}))$ et
$f'_*(g^{\prime*}(\mathcal{F}))$ sont tous deux égaux au
$\mathcal{O}_{S'}$-Module associé au $A'$-module $A'\otimes_A M$ (où $M$ est
considéré comme $A$-module), et que $u$ est l'homomorphisme associé à
l'identité (\hyperref[I.1.6.3-fr]{I, 1.6.3},
\hyperref[I.1.6.5-fr]{1.6.5} et \hyperref[I.1.6.7-fr]{1.6.7}).

\begin{remark}[1.5.3]
\label{II.1.5.3-fr}
On se gardera de croire que (\hyperref[II.1.5.2-fr]{1.5.2}) reste valable
lorsque $X$ n'est plus supposé affine au-dessus de $S$, même lorsque
$S'=\operatorname{Spec}(k(s))$ ($s\in S$) et que $S'\to S$ est le morphisme
canonique (\hyperref[I.2.4.5-fr]{I, 2.4.5}) --- auquel cas $X'$ n'est autre
que la \emph{fibre} $f^{-1}(s)$ (\hyperref[I.3.6.2-fr]{I, 3.6.2}). En
d'autres termes, lorsque $X$ n'est pas affine au-dessus de $S$, l'opération
